\chapter{Belief Fibers and Exponential Families}
\label{ch:expfamily}

The objects carried by the belief and model associated bundles of
\Cref{ch:geometry} are normalized laws, while generative fibers carry kernels.
Coordinates, smoothness, domination, and exponential-family structure are
additional choices. This chapter states that hierarchy before selecting the
finite-dimensional exponential subclass used by the general coarse theory.

\section{Law fibers before coordinates}
\label{sec:exp-general-fibers}

\definitionheading{Law and kernel fibers}{def:exp-law-kernel-fibers}
For a standard Borel space \((\mathsf Z,\mathscr Z)\), let
\begin{equation}
\mathcal B(\mathsf Z)\subseteq\mathcal P(\mathsf Z)
\label{eq:exp-law-fiber}
\end{equation}
be a declared collection of normalized probability laws. For standard Borel
spaces \(\mathsf X,\mathsf Y\), let
\begin{equation}
\mathcal K(\mathsf X,\mathsf Y)
\subseteq\operatorname{Markov}(\mathsf X,\mathsf Y)
\label{eq:exp-kernel-fiber}
\end{equation}
be a declared collection of normalized Markov kernels. A law is not a kernel,
and neither declaration supplies a common density, affine chart, manifold,
Fisher metric, or convex cone. \status{DEFINITION}

The belief and model bundles may use different law fibers. Cross-bundle maps
\(\Phi:\mathcal E_b\to\mathcal E_m\) and
\(\widetilde\Phi:\mathcal E_m\to\mathcal E_b\) act between those associated
bundles and are not iterations on either fiber. To prevent collision with
that reserved notation, a generic endomorphism of a parameter set will be
written \(\Psi\).

\definitionheading{Partial realization}{def:exp-partial-realization}
A parametrized model consists of a parameter space \(\Theta\), a domain
\(\Theta_{\rm law}\subseteq\Theta\), and a map
\begin{equation}
\mathfrak r:\Theta_{\rm law}\longrightarrow\mathcal B(\mathsf Z),
\qquad
\vartheta\longmapsto Q_\vartheta .
\label{eq:exp-realization-map}
\end{equation}
Outside \(\Theta_{\rm law}\), a parameter may still denote an algebraic
operator, but it does not denote a law through \(\mathfrak r\).
\status{DEFINITION}

A bimeasurable gauge action \(g:\mathsf Z\to\mathsf Z\) acts on laws by
\begin{equation}
Q\longmapsto g_\#Q ,
\label{eq:exp-general-pushforward}
\end{equation}
and on kernels by
\begin{equation}
(g_{\rm out}\cdot K\cdot g_{\rm in}^{-1})(x,B)
=K(g_{\rm in}^{-1}x,g_{\rm out}^{-1}B).
\label{eq:exp-kernel-action}
\end{equation}
These formulas type the action. \status{DEFINITION}

Closure of a declared fiber under the action is a separate hypothesis.
\status{HYPOTHESIS}

\section{Dominated and smooth tiers}
\label{sec:exp-dominated-nondominated}

Domination is not automatic. An uncountable family of Dirac laws on an
uncountable standard Borel space admits no common sigma-finite dominating
measure: any such measure would have to assign positive mass to uncountably
many disjoint atoms, which a sigma-finite measure cannot do.
\status{ESTABLISHED}

Hence density charts cannot be inferred from the general law fiber.
\status{ESTABLISHED}

The differential tier used below is a finite-dimensional smooth
parametrized-measure model. Differentiability in quadratic mean,
square-integrable scores, and the interchange of differentiation and
integration are hypotheses. The Fisher form is a Riemannian metric only after
nondegeneracy is proved. Parametrized-measure models supply an intrinsic
framework for this tier \citep{AyJostLeSchwachhoefer2018}.
\status{HYPOTHESIS}

Stratified and infinite-dimensional tiers require their own tangent and
integration theories; no universal global Fisher metric or natural chart is
claimed.
\status{NOT-CLAIMED}

\section{Regular finite-dimensional exponential families}
\label{sec:exp-family}

\definitionheading{Canonical family}{def:exp-canonical-family}
Let \(\nu\) be a nonzero sigma-finite measure on
\((\mathsf Z,\mathscr Z)\), let \(T:\mathsf Z\to\mathsf T\) be measurable
into a finite-dimensional inner-product space, and define
\begin{equation}
\mathsf A(\theta)
=\log\int_{\mathsf Z}e^{\langle\theta,T(z)\rangle}\nu(dz),
\qquad
\mathsf N=\{\theta:\mathsf A(\theta)<\infty\}.
\label{eq:exp-logpartition}
\end{equation}
For \(\theta\in\mathsf N\),
\begin{equation}
q_\theta(z)
=\exp\{\langle\theta,T(z)\rangle-\mathsf A(\theta)\},
\qquad
Q_\theta(dz)=q_\theta(z)\nu(dz).
\label{eq:exp-density}
\end{equation}
Nothing is asserted for \(\theta\notin\mathsf N\). \status{DEFINITION}

\propositionheading{Finite-domain interface to ELBO coordinates}{prop:exp-finite-domain-elbo-interface}
The map $(\theta,z)\mapsto q_\theta(z)$ in \eqref{eq:exp-density} is jointly measurable on $\mathsf N\times\mathsf Z$, and every law, density ratio, Fisher tensor, or ELBO coordinate built from this canonical representation is restricted to parameters in $\mathsf N$. For a comparison of several canonical representations, either the comparison is made intrinsically at the measure level or the representations must have one declared common sigma-finite base measure; its density domain is the intersection of their natural domains. An extended-ELBO update additionally requires a positive-finite selected evidence slice at every compared parameter, while the classical expected-log-density representation separately requires \Cref{hyp:prob-classical-split-elbo}. An algebraic operator that leaves $\mathsf N$ therefore ceases to define a probability or ELBO coordinate even if it remains an element of the operator cone. \status{ESTABLISHED}

\paragraph{Proof.}
Lower semicontinuity makes $\mathsf A$ Borel measurable, and measurability of $T$ makes $(\theta,z)\mapsto\langle\theta,T(z)\rangle$ jointly measurable; exponentiation proves the first claim. Normalization is finite exactly on $\mathsf N$ by \eqref{eq:exp-logpartition}. The remaining statements are direct domain checks against \Cref{def:exp-partial-realization,def:elbo-extended,hyp:prob-classical-split-elbo}: a density comparison needs a fixed reference, the extended functional needs a normalized posterior slice with finite positive mass, and its classical split needs the additional support and log-integrability hypotheses. \(\square\)

\propositionheading{Convexity and lower semicontinuity}{prop:exp-logpartition-convex-lsc}
\(\mathsf A\) is convex and lower semicontinuous as an extended-real
function, and \(\mathsf N\) is convex. \status{ESTABLISHED}

\paragraph{Proof.}
For \(0<t<1\), H\"older's inequality gives
\begin{equation}
\int e^{\langle t\theta+(1-t)\vartheta,T\rangle}d\nu
\leq
\left(\int e^{\langle\theta,T\rangle}d\nu\right)^t
\left(\int e^{\langle\vartheta,T\rangle}d\nu\right)^{1-t},
\label{eq:exp-holder}
\end{equation}
and taking logarithms proves convexity and convexity of the effective domain.
Fatou's lemma applied to a convergent parameter sequence proves lower
semicontinuity. \(\square\)

\propositionheading{Cumulants and Fisher information}{prop:exp-cumulants-fisher}
At every interior parameter where local exponential domination permits
differentiation under the integral,
\begin{equation}
\nabla\mathsf A(\theta)=\E_\theta[T],
\qquad
\nabla^2\mathsf A(\theta)=\operatorname{Cov}_\theta(T).
\label{eq:exp-cumulants}
\end{equation}
The Hessian is the Fisher information in the natural chart.
\status{ESTABLISHED}

\paragraph{Proof.}
Differentiate the normalizer once and twice. The score is
\(T-\E_\theta T\), so its second moment is the covariance in
\eqref{eq:exp-cumulants}. Local exponential domination follows by enclosing
the parameter in a compact subset of the interior and applying the elementary
bound
\begin{equation}
|x|^r\leq C_{r,\varepsilon}
\left(e^{\varepsilon x}+e^{-\varepsilon x}\right)
\label{eq:exp-domination}
\end{equation}
to each coordinate of \(T\). \(\square\)

\hypothesisheading{Regularity and minimality}{hyp:exp-regular-minimal}
The family is regular, so \(\mathsf N\) is open, and minimal, so no nonzero
\(a\in\mathsf T\) makes \(\langle a,T\rangle\) constant
\(\nu\)-almost everywhere. Regularity excludes natural-domain boundary
points; minimality excludes affine redundancy. \status{HYPOTHESIS}

\propositionheading{Minimality and nondegenerate information}{prop:exp-minimal-fisher-nondegenerate}
Under the preceding hypothesis, \(\nabla^2\mathsf A(\theta)\) is positive
definite, \(\mathsf A\) is strictly convex, and
\(\theta\mapsto Q_\theta\) is injective. Failure of minimality makes the
Hessian singular. \status{ESTABLISHED}

\paragraph{Proof.}
For every \(a\),
\[
a^\top\nabla^2\mathsf A(\theta)a
=\operatorname{Var}_{Q_\theta}(\langle a,T\rangle).
\]
The variance vanishes exactly when the displayed statistic is almost surely
constant. Since \(q_\theta>0\), \(Q_\theta\) and \(\nu\) have the same null
sets, so minimality gives positive definiteness. Strict convexity follows.
Equality of two densities makes
\(\langle\theta-\vartheta,T\rangle\) constant almost everywhere, and
minimality then gives \(\theta=\vartheta\). \(\square\)

\section{Operator and probability layers}
\label{sec:exp-layers}

\definitionheading{Operator layer}{def:exp-operator-layer}
Let \(\mathcal C\subseteq\mathsf T\) be a separately declared closed convex
cone preserved by the algebraic maps under study. Its probability layer is
\(\mathcal C\cap\mathsf N\). The cone is chosen for algebraic closure, not
normalization, and no such cone is asserted for an arbitrary law fiber.
\status{DEFINITION}

\propositionheading{A natural-domain boundary carries no law}{prop:exp-domain-boundary-no-law}
Under regularity, if \(\theta^\star\in\partial\mathsf N\), then
\begin{equation}
\int e^{\langle\theta^\star,T(z)\rangle}\nu(dz)=+\infty.
\label{eq:exp-boundary-diverges}
\end{equation}
For \(\theta_\tau=(1-\tau)\theta_0+\tau\theta^\star\), with
\(\theta_0\in\mathsf N\) and \(w=\theta^\star-\theta_0\),
\begin{equation}
\mathsf A(\theta_\tau)\to+\infty,\qquad
\langle w,\nabla\mathsf A(\theta_\tau)\rangle\to+\infty,\qquad
\int_0^1\operatorname{Var}_{Q_{\theta_\tau}}(\langle w,T\rangle)d\tau
=+\infty.
\label{eq:exp-boundary-blowup}
\end{equation}
\status{ESTABLISHED}

\paragraph{Proof.}
Regularity puts \(\theta^\star\) outside \(\mathsf N\), proving
\eqref{eq:exp-boundary-diverges}. Convexity puts every
\(\theta_\tau\), \(\tau<1\), in \(\mathsf N\), while lower
semicontinuity forces the first limit. For
\(g(\tau)=\mathsf A(\theta_\tau)\), convexity makes \(g'\)
nondecreasing. If \(g'\) were bounded above, then \(g\) would remain bounded,
a contradiction. The identities
\[
g'=\langle w,\nabla\mathsf A\rangle,\qquad
g''=\operatorname{Var}_{Q_{\theta_\tau}}(\langle w,T\rangle)
\]
and integration of \(g''\) prove the remaining assertions. \(\square\)

Thus the declared representation supplies no finite log partition, mean
parameter, Fisher metric, relative entropy from the nonexistent endpoint
law, or instance of the exact ELBO at \(\theta^\star\).
\status{ESTABLISHED}

This does not rule out a quotient, pin, new reference measure, or boundary
stratum; each is a different declared object. \status{NOT-CLAIMED}

\propositionheading{An operator fixed point need not be a law}{prop:exp-fixed-point-no-law}
Let \(\Psi:\mathcal C\to\mathcal C\) be continuous. Every iterate may lie in
\(\mathcal C\cap\mathsf N\) while the unique fixed point lies on
\(\partial\mathsf N\). For
\(\mathsf Z=\R\), \(T(z)=-z^2/2\), and Lebesgue \(\nu\),
\begin{equation}
\mathsf A(\theta)=-\tfrac12\log\theta+\tfrac12\log2\pi,
\qquad
\mathsf N=(0,\infty),
\qquad
\mathcal C=[0,\infty),
\label{eq:exp-scale-witness}
\end{equation}
the map \(\Psi(\theta)=\theta/2\) has this property.
\status{ESTABLISHED}

\paragraph{Proof.}
Starting at one gives \(\theta_\ell=2^{-\ell}>0\), while
\(\theta_\ell\to0\), the unique fixed point, and
\(\mathsf A(\theta_\ell)\to+\infty\). The constant kernel at zero has
infinite Lebesgue integral. \(\square\)

\section{Declared repairs}
\label{sec:exp-repairs}

\definitionheading{Pinning}{def:exp-pinning}
Choose \(\theta_{\rm m}\in\mathsf N\) and replace
\(\mathcal C\) by \(\theta_{\rm m}+\mathcal C\); the pin does not flow.
\status{DEFINITION}

\propositionheading{Pinning under a recession condition}{prop:exp-pinning-recession}
If \(\theta_{\rm m}\in\operatorname{int}\mathsf N\) and every
\(\eta\in\mathcal C\) is a recession direction of \(\mathsf N\), then
\(\theta_{\rm m}+\mathcal C\subseteq\operatorname{int}\mathsf N\).
\status{ESTABLISHED}

\paragraph{Proof.}
Translate an open ball around \(\theta_{\rm m}\) by \(\eta\). The recession
hypothesis keeps the translated ball inside \(\mathsf N\), so its center is
interior. \(\square\)

The pin changes the family, may change the fixed set, and introduces a
parameter whose removal need not commute with the long-scale limit. Whether
the mass sector is relevant, marginal, or irrelevant is an open spectral
question for the declared scale map. \status{OPEN}

\definitionheading{Quotienting}{def:exp-quotienting}
Suppose \(\mathsf Z=\R^n\), \(\nu\) is Lebesgue measure, and the energy is
invariant along a nonzero linear subspace \(\mathcal V\). Declare the new
sample space \(\mathcal W=\mathcal V^\perp\), its Lebesgue measure, and the
restricted statistic. \status{DEFINITION}

\propositionheading{Quotient normalization is not full-space normalization}{prop:exp-quotient-normalization}
Writing \(z=v+w\) in the orthogonal splitting,
\begin{equation}
\int_{\R^n}e^{\langle\theta,T(z)\rangle}dz
=\lambda_{\mathcal V}(\mathcal V)
\int_{\mathcal W}e^{\langle\theta,T(w)\rangle}dw
=+\infty.
\label{eq:exp-quotient-split}
\end{equation}
The quotient law is normalized exactly when the second factor is finite.
\status{ESTABLISHED}

\paragraph{Proof.}
The change of variables has unit Jacobian and the integrand is independent of
\(v\). Tonelli's theorem gives \eqref{eq:exp-quotient-split}; the first factor
is infinite. \(\square\)

\definitionheading{Retaining a mass sector}{def:exp-mass-sector}
Enlarge the state by a declared coercive parameter that remains in the
natural domain under the scale map. This is neither pinning nor quotienting:
the mass is a running coordinate and its scaling law must be stated.
\status{DEFINITION}

A Moore--Penrose inverse and pseudodeterminant inserted into an interior
full-space density formula do not repair normalization. They compute the
finite quotient expression, whose dimension and reference measure differ,
while the original integral remains infinite by
\eqref{eq:exp-quotient-split}. \status{ESTABLISHED}

\section{Relative entropy and affine information projection}
\label{sec:exp-information-projection}

For a regular minimal family, let
\(\tau_\theta=\nabla\mathsf A(\theta)\) and let
\(\mathsf A^\ast\) be the convex conjugate. Direct expectation of the
log-density ratio gives
\begin{equation}
\KL(Q_\vartheta\Vert Q_\theta)
=\mathsf A(\theta)-\mathsf A(\vartheta)
-\langle\nabla\mathsf A(\vartheta),\theta-\vartheta\rangle
=D_{\mathsf A}(\theta,\vartheta),
\label{eq:exp-kl-bregman}
\end{equation}
and Legendre duality gives
\begin{equation}
\KL(Q_\vartheta\Vert Q_\theta)
=D_{\mathsf A^\ast}(\tau_\vartheta,\tau_\theta).
\label{eq:exp-kl-dual-bregman}
\end{equation}
This is the exact exponential-family divergence; in general it is not a
finite-displacement Fisher quadratic
\citep{Csiszar1967,Brown1986,Amari2016}. \status{ESTABLISHED}

\propositionheading{Affine mean constraints attain a unique projection}{prop:exp-affine-mean-projection}
Assume the extended \(\mathsf A\) is Legendre, as in a regular minimal
canonical family, and let
\(\mathsf M=\operatorname{int}\operatorname{dom}\mathsf A^\ast\).
If \(C\) is a closed affine subspace with
\(C\cap\mathsf M\neq\varnothing\), then for every
\(\tau_\theta\in\mathsf M\) there is a unique
\(\widehat\tau\in C\cap\mathsf M\) minimizing
\begin{equation}
\tau\longmapsto
D_{\mathsf A^\ast}(\tau,\tau_\theta).
\label{eq:exp-mean-projection}
\end{equation}
The corresponding natural parameter is uniquely
\(\widehat\vartheta=\nabla\mathsf A^\ast(\widehat\tau)\).
\status{ESTABLISHED}

\paragraph{Proof.}
The finite-dimensional Legendre Bregman projection theorem gives existence
\citep{Bauschke2003},
uniqueness, and interiority on every nonempty closed convex feasible set
meeting the interior domain. Strict convexity gives uniqueness, and essential
smoothness excludes boundary attainment. For
\(C=\{\tau:B\tau=b\}\), the solution equivalently satisfies
\begin{equation}
B\widehat\tau=b,\qquad
\nabla\mathsf A^\ast(\widehat\tau)-\theta+B^\top\lambda=0,
\label{eq:exp-projection-kkt}
\end{equation}
which also proves the stated natural-parameter formula. \(\square\)

Feasibility, closedness, and convexity are load bearing. For a Bernoulli mean
domain \((0,1)\), \(C=\{0\}\) is infeasible; an open feasible interval may
have a finite unattained infimum; and a closed nonconvex two-point set can
have two minimizers. Gauge compatibility and commutation of this projection
with a coarse parameter map remain open and require an intertwining theorem;
existence alone does not provide them. \status{OPEN}

\section{Boundary of this chapter}
\label{sec:exp-downstream-assumptions}

This chapter supplies the general law-fiber hierarchy, the exponential-family
interior, its operator boundary, and the Bregman projection theorem. It does
not contain graph aggregation, Gaussian determinants, Schur complements, or
matrix pencils. Graph-exponential energy closure is developed in
\Cref{ch:coarsegraining}; the multivariate-Gaussian realization and its
quadratic boundary rates appear later. An operator RG additionally requires
a preserved parameter category, scale comparisons, rescaling, and topology.
A law-level RG additionally requires normalized laws or kernels, compatible
reference measures, and finite normalizers at every scale. A Fisher or ELBO
reading additionally requires the hypotheses of the corresponding earlier
chapters. \status{HYPOTHESIS}
